\documentclass[12pt]{article}
\usepackage[utf8]{inputenc}
\usepackage[T1]{fontenc}
\usepackage{hyperref}
\title{Fiwaga Language Resources for K-12 Curriculum Development}
\author{Compiled by Claude Sonnet 5.5 (claude-sonnet-5-5) for Andr\'as Kornai}
\date{}
\begin{document}
\maketitle

\section{Introduction}
Fiwaga (also Fimaga, Fiwage; ISO 639-3 fiw) is a Papuan language of Southern Highlands Province, Papua New Guinea, near Lake Kutubu. It forms the East Kutubuan family with Foi. Speaker counts in the sources conflict (about 300 native speakers in a 1981 count versus roughly 1,000 people in the ethnic group). Very little exists for this language; this report is based on a few searches.

\section{Primary Educational Resources}
None found. Joshua Project states that the written form is essentially undeveloped.

\section{Linguistic and Reference Materials}
\begin{itemize}
\item Wikipedia, Fiwaga language: \url{https://en.wikipedia.org/wiki/Fiwaga_language}
\item Wikipedia, East Kutubuan languages (Fiwaga and Foi): \url{https://en.wikipedia.org/wiki/East_Kutubuan_languages}
\item New Guinea World, Lake Kutubu language group overview: \url{https://newguineaworld.linguistik.uzh.ch/families/papuan-gulf/kikori-river/lake-kutubu/start}
\end{itemize}

\section{Religious and Community Resources}
Joshua Project reports that Bible translation has begun but is far from complete, with no audio Bible or Jesus Film:
\url{https://joshuaproject.net/languages/fiw} and \url{https://joshuaproject.net/people_groups/11215/PP}

\section{Digital and Audio Resources}
None found.

\section{Curriculum Development Context}
There are no known textbooks, primers, dictionaries or audio in Fiwaga. A curriculum would have to begin with documentation, orthography development and community consultation. The closely related Foi language may have more material, but I did not search for it and do not cite any.

\end{document}